\clearpage\pagestyle{otMain}\fontsize{11}{15}\selectfont\corrigesChapitre{7}{Grandeurs quotients et conversions}
\begin{blocLivre}{cor-7-1}\begin{corrige}{1}{Des longueurs et des aires}\demarche{Distinguer le facteur des longueurs de celui des aires.}
On obtient $\res{3800\text{ m}}$, $\res{4200\text{ cm}^2}$ et $\res{0{,}65\text{ m}^2}$. Les deux dernières conversions utilisent le facteur 10000.
\end{corrige}\end{blocLivre}
\begin{blocLivre}{cor-7-2}\begin{corrige}{2}{Volumes et contenances}\demarche{Utiliser les équivalences entre volumes et contenances.}
$2{,}7\text{ m}^3=\res{2700\text{ L}}$. Comme un mL est un cm$^3$, 450 mL valent $\res{450\text{ cm}^3}$, soit $\res{0{,}45\text{ L}}$.
\end{corrige}\end{blocLivre}
\begin{blocLivre}{cor-7-3}\begin{corrige}{3}{Une durée décimale}\demarche{Convertir la partie d’heure avec 60.}
$0{,}75\times60=45$, donc $\res{1\text{ h }45\text{ min}}$. Inversement, $42/60=0{,}7$, donc $\res{2{,}7\text{ h}}$.
\end{corrige}\end{blocLivre}
\begin{blocLivre}{cor-7-4}\begin{corrige}{4}{Une vitesse moyenne}\demarche{Mettre toute la durée dans l’unité choisie.}
La durée vaut 2,5 h. La vitesse moyenne est $150/2{,}5=\res{60\text{ km/h}}$. Diviser par 2,30 aurait mal traduit les minutes.
\end{corrige}\end{blocLivre}
\begin{blocLivre}{cor-7-5}\begin{corrige}{5}{Un débit}\demarche{Calculer le volume d’une minute.}
Le débit est $84/7=\res{12\text{ L/min}}$. Le volume en 25 minutes vaut $12\times25=\res{300\text{ L}}$.
\end{corrige}\end{blocLivre}
\begin{blocLivre}{cor-7-6}\begin{corrige}{6}{Convertir une vitesse}\demarche{Revenir aux unités de base pour justifier le facteur 3,6.}
$90/3{,}6=\res{25\text{ m/s}}$. Dans l’autre sens, $12\times3{,}6=\res{43{,}2\text{ km/h}}$. Un contrôle direct emploie 1000 m et 3600 s.
\end{corrige}\end{blocLivre}
\begin{blocLivre}{cor-7-7}\begin{corrige}{7}{Le temps de réaction}\demarche{Séparer la réaction du freinage.}
La vitesse vaut $54/3{,}6=15$ m/s. La distance pendant 0,8 s est $15\times0{,}8=\res{12\text{ m}}$. Ce calcul ne comprend pas la distance de freinage.
\end{corrige}\end{blocLivre}
\begin{blocLivre}{cor-7-8}\begin{corrige}{8}{Une masse volumique}\demarche{Multiplier le volume par la masse d’une unité de volume.}
La masse vaut $\rho V=2{,}5\times360=\res{900\text{ g}}$, soit $\res{0{,}9\text{ kg}}$. Les cm$^3$ du volume sont cohérents avec ceux du dénominateur de la masse volumique.
\end{corrige}\end{blocLivre}
\begin{blocLivre}{cor-7-9}\begin{corrige}{9}{Une consommation électrique}\demarche{Convertir la durée avant de calculer l’énergie.}
45 minutes valent 0,75 h. L’énergie vaut $1{,}2\times0{,}75=\res{0{,}9\text{ kWh}}$. Le coût vaut $0{,}9\times0{,}24=\res{0{,}216\text{ €}}$, soit 0,22 € au centime. Le tarif est une donnée de l’exercice.
\end{corrige}\end{blocLivre}
\begin{blocLivre}{cor-7-10}\begin{corrige}{10}{Un téléchargement}\demarche{Remplacer les préfixes par des puissances de 10 avant de diviser.}
$2{,}4\text{ Go}=2{,}4\times10^9$ octets et $40\text{ Mo/s}=4\times10^7$ octets par seconde. La durée vaut $2{,}4\times10^9/(4\times10^7)=0{,}6\times10^2=\res{60\text{ s}}$, soit une minute.
\end{corrige}\end{blocLivre}
\begin{blocLivre}{cor-7-11}\begin{corrige}{11}{Aller et retour}\demarche{Additionner distances et durées avant de diviser.}
L’aller dure $12/18=2/3$ h ; le retour dure 1 h. Le total est 24 km en $5/3$ h. La vitesse moyenne vaut $24/(5/3)=\res{14{,}4\text{ km/h}}$. La moyenne 15 km/h des deux vitesses ne convient pas, car les durées diffèrent.
\end{corrige}\end{blocLivre}
\begin{blocLivre}{cor-7-12}\begin{corrige}{12}{Deux robinets}\demarche{Additionner les fractions de cuve remplies par minute.}
En une minute, ils remplissent respectivement $1/30$ et $1/45$ de la cuve. Ensemble : $1/30+1/45=1/18$. Le remplissage prend donc $\res{18\text{ min}}$. Additionner les deux durées aurait ignoré l’addition des débits.
\end{corrige}\end{blocLivre}
\begin{blocLivre}{cor-7-13}\begin{corrige}{13}{Choisir assez de peinture}\demarche{Calculer l’aire utile, les couches, puis l’achat entier.}
L’aire à peindre est $5\times2{,}6-2=11$ m$^2$. Les deux couches représentent 22 m$^2$ à couvrir, donc $22/8=2{,}75$ L. Un pot ne suffit pas ; il faut $\res{2\text{ pots}}$. On n’arrondit pas le besoin au-dessous.
\end{corrige}\end{blocLivre}
\begin{blocLivre}{cor-7-14}\begin{corrige}{14}{Une fuite discrète}\demarche{Construire le débit puis multiplier par la durée totale.}
Le débit est 0,1 mL/s. Une journée compte 86400 s, d’où $8640$ mL, soit $\res{8{,}64\text{ L}}$. En 30 jours : $\res{259{,}2\text{ L}}$. Le résultat dépend du volume de goutte choisi dans le modèle.
\end{corrige}\end{blocLivre}
\begin{blocLivre}{cor-7-15}\begin{corrige}{15}{Une livraison de gravier}\demarche{Employer les mètres pour les trois dimensions du volume.}
On convertit l’épaisseur en 0,06 m. Le volume vaut $12\times1{,}5\times0{,}06=1{,}08$ m$^3$. La masse est $1{,}08\times1{,}6=1{,}728$ tonne. Elle dépasse 1,5 tonne : $\res{2\text{ voyages}}$ suffisent et un seul est impossible.
\end{corrige}\end{blocLivre}
